\section{系统总体架构与分层结构}

本节阐述了全尺寸金融市场仿真平台的总体架构和分层结构。整个平台围绕“\textbf{核心仿真层}、\textbf{服务支撑层}、\textbf{交互与监控层}”三大层次构建，每一层负责不同的功能模块，并通过明确定义的接口协作。采用这种分层设计有助于降低各组件之间的耦合度，便于系统的开发、测试和后期扩展。

\subsection{总体架构概述}

总体架构如图~\ref{fig:component_overview} 所示。核心仿真层位于底部，负责实现市场行为的核心逻辑，包括事件调度、消息传递、订单簿管理和撮合引擎。服务支撑层构建在核心仿真层之上，提供数据管理、代理框架、自校准、实验配置等辅助功能，既为仿真提供必要的支撑，又使用户能够灵活配置和扩展仿真实验。最上层的交互与监控层面向最终用户，通过程序化接口、命令行和可视化界面提供交互入口，并负责展示实时指标、日志和分析结果。

\subsection{核心仿真层}

核心仿真层是平台的运行基础，实现仿真的时间推进和市场微观行为。这一层主要包含以下模块：

\begin{itemize}
  \item \textbf{事件调度器（EventScheduler）}：维护全局虚拟时间和优先级事件队列，按时间顺序调度并执行各类事件。为支持大规模并行，可将不同资产或代理划分到独立的线程或协程中，并在时间窗口边界进行同步，确保全局时间一致。
  \item \textbf{消息总线（MessageBus）}：采用消息驱动模式实现组件之间的解耦，负责不同模块与代理之间的消息发布与订阅，包括订单提交、成交报告、行情更新等。消息格式统一定义，便于代理扩展。
  \item \textbf{市场环境（Market Environment）}：模拟真实交易所，包括多资产订单簿、撮合引擎和拍卖机制。订单簿使用跳表或红黑树等高效数据结构，实现价格–时间优先的插入、删除和查询。撮合引擎支持限价单、市场单、隐藏单等多种订单类型，并在连续竞价与集合竞价间切换。跨资产协调器确保多资产行情和订单同步，处理跨市场延迟和套利行为。
\end{itemize}

通过这一层的抽象设计，底层仿真逻辑与上层功能解耦，使得市场机制的改进（如加入新的订单类型或撮合规则）不影响服务层和交互层的实现。

\subsection{服务支撑层}

服务支撑层为核心仿真提供通用功能和扩展能力，主要由以下模块组成：

\begin{itemize}
  \item \textbf{代理框架（Agent Framework）}：定义统一的代理基类和生命周期，提供市场观察、下单、撤单等基本接口。内置多种代理模型（零智商代理、做市商、统计套利代理、强化学习代理等），并允许用户继承基类实现自定义策略。代理框架屏蔽了内核细节，使策略开发者能够专注于决策逻辑。
  \item \textbf{自校准系统（Calibration System）}：负责监测市场统计指标并动态调整背景代理参数，使仿真输出与目标市场特征保持一致。通过指标监控器收集订单簿深度、成交量、收益率分布等数据，与预设目标比较后，调用参数控制器调节到达率、订单量分布或价格偏好，实现在线校准。
  \item \textbf{数据管理模块（Data Management）}：提供统一的数据接口，支持读取历史行情、生成合成数据及访问实时行情。采用列式存储与缓存机制提升数据访问效率，并提供异步查询接口供代理拉取数据。
  \item \textbf{配置与实验管理（Config\,&\,Experiment Manager）}：解析用户提供的配置文件或命令行参数，初始化仿真环境，管理实验的生命周期，并支持参数扫描和批量实验调度。
  \item \textbf{监控与日志（Monitoring\,&\,Logging）}：收集并存储仿真过程中产生的事件、性能指标和代理状态，为自校准系统和用户分析提供数据来源。该模块还负责输出日志、生成报告，支持外部监控工具接入。
\end{itemize}

服务支撑层既向下依赖核心仿真提供基础运行环境，又向上为交互与监控层提供数据和配置支持，是系统灵活性和可扩展性的关键所在。

\subsection{交互与监控层}

交互与监控层面向最终用户，提供访问接口和实时反馈。主要组件包括：

\begin{itemize}
  \item \textbf{程序化 API}：通过 HTTP/gRPC 等方式对外暴露接口，使外部程序（如强化学习算法、回测框架）能够提交订单、查询市场状态、注入策略等。接口定义遵循 RESTful 设计，支持异步调用。
  \item \textbf{命令行和脚本接口}：为实验人员提供便捷的命令行工具，支持启动仿真、设置参数、插入策略、查看仿真状态等功能，便于批量实验和自动化运行。
  \item \textbf{可视化界面}：提供图形化监控界面，展示实时行情、订单簿深度、成交记录以及各代理的持仓和盈亏等信息。界面支持交互操作，如实时调整策略参数或政策设置。
  \item \textbf{日志与分析}：将监控模块收集的数据经过分析和可视化，形成用户友好的报告；支持导出仿真记录和统计结果。
\end{itemize}

这一层与服务支撑层通过清晰的 API 交互，既可以将复杂的仿真能力以服务形式提供给外部系统，又能为用户提供直观的使用体验和实时反馈。

\subsection{组件关系图}

图~\ref{fig:component_overview} 为平台的总体组件关系图。在绘图中，按照三层架构排列各模块，通过箭头指示依赖关系和信息流向。该图在附录中以 TikZ 形式绘制，并可供后续设计和实现参考。

\begin{figure}[htbp]
  \centering
  \resizebox{0.95\linewidth}{!}{\input{figures/component_overview.tex}}
  \caption{系统总体组件关系图。图中按照分层结构绘制主要模块及其连接关系。}
  \label{fig:component_overview}
\end{figure}